\section{Introducción}
En el paradigma de gestión empresarial contemporáneo, las organizaciones operan en entornos altamente dinámicos donde la capacidad de respuesta y la precisión estratégica dependen de la calidad con la que se procesan y analizan los datos corporativos. Tradicionalmente, los sistemas transaccionales en línea (OLTP, Online Transaction Processing) han sido diseñados para optimizar la inserción, actualización y consulta puntual de registros operativos diarios con un alto grado de normalización formal (tercera forma normal o 3NF). Si bien esta arquitectura relacional garantiza la integridad referencial y la concurrencia en tareas como la facturación o la gestión de inventarios, resulta ineficiente e inviable cuando se enfrenta a cargas de trabajo analíticas complejas que involucran agregaciones masivas, análisis de tendencias históricas y minería de datos.La ejecución de consultas analíticas pesadas directamente sobre bases de datos de producción suele provocar contención de recursos, bloqueos de tablas (table locks) y degradación severa del tiempo de respuesta del servicio transaccional. Para mitigar esta fricción arquitectónica, surge la disciplina del Almacenamiento de Datos (Data Warehousing), cuyo propósito es desacoplar las operaciones transaccionales del análisis estratégico corporativo mediante la creación de un repositorio centralizado, integrado y no volátil: la Bodega de Datos (Data Warehouse).Bajo la definición clásica de W. H. Inmon y los principios de modelado dimensional popularizados por Ralph Kimball, un Data Warehouse no es una simple réplica de la base de datos operativa, sino un ecosistema estructurado específicamente para el Procesamiento Analítico en Línea (OLAP, Online Analytical Processing). La construcción de este repositorio implica estructurar los datos bajo una "única fuente de verdad" (Single Source of Truth), resolviendo las inconsistencias semánticas entre subsistemas y permitiendo la navegación multidimensional de los indicadores de rendimiento (KPI). En el centro de esta disciplina se encuentra el Modelo Dimensional, que organiza los datos en tablas de hechos (que capturan métricas cuantitativas del negocio) y tablas de dimensiones (que aportan el contexto de análisis: quién, cuándo, dónde y qué).   Para la empresa Spring Valley, el reto de consolidar sus operaciones comerciales responde a la imperiosa necesidad de transformar registros aislados de ventas, cobros y catálogo de productos en un activo analítico de alto valor. Construir una bodega de datos para el área comercial implica diseñar e implementar flujos de extracción, transformación y carga (ETL), definir granularidades atómicas de venta que preserven el detalle transaccional de facturación y parametrizar un esquema en estrella (Star Schema) optimizado para motores relacionales modernos.El presente documento técnico establece la guía arquitectónica y metodológica para el despliegue del Data Warehouse de Spring Valley. A lo largo de sus secciones se detallan el análisis de requerimientos técnicos y de hardware/software, el cronograma de ejecución bajo el ciclo de vida de Kimball, el modelado dimensional relacional (con el diccionario de datos y tipos de datos SQL correspondientes), la persistencia e infraestructura contenerizada en Docker con PostgreSQL 18-alpine, y las estrategias de ingeniería implementadas para garantizar la gobernanza, escalabilidad y rendimiento en la toma de decisiones ejecutivas.